\section{Research design}\label{sec:design}

The comparison at the centre of this report is between an expiry Thursday and an ordinary
session. Two things stand between that comparison and a conclusion about the settlement
mechanism, and the design is built around both.

The first is the market environment. A given month's volatility, index level and news flow
affect the expiry session and everything around it, so the comparison is made within the
month: each of the \DesignExpirySessions{} expiry Thursdays in 2022 is paired with a control
session in the same calendar month, giving \DesignSessions{} sessions. Every test is paired at
the level of a security and a month, so a security that is simply busier than another
contributes nothing to any comparison.

The second is the calendar. An expiry Thursday is also the last Thursday of the month, when
index rebalancing, month-end portfolio flows and index option expiry all fall due. Matching
within the month does nothing about this: it is a property of the day itself. A study that
stops here can report that the settlement window is different and cannot say why, which is
the position the first version of this project was in.

\subsection{Three groups}

The answer is to include securities for which the settlement mechanism does not exist.

\paragraph{Liquid and illiquid, both with derivatives.} Both groups carry an expiring
contract that settles against their own cash VWAP, so both face the same incentive. What
differs is the cost of acting on it, which Section~\ref{sec:results} quantifies directly by
sweeping the displayed book. If the settlement incentive is being acted on, the effect should
be larger where acting on it is cheaper.

This only works if the groups actually differ in liquidity. The first version of this study
used five ``liquid'' and five ``illiquid'' securities, every one of which was a large-cap
index constituent; the contrast measured nothing, because both groups were among the most
traded securities on the exchange.

\paragraph{Placebo, without derivatives.} No contract settles against these securities, so
nobody has any reason to care about their closing half hour for settlement purposes. They are
still subject to everything else that happens on the last Thursday of a month. If spreads
widen and cancellations rise there too, the effect belongs to the calendar; if they do not,
it is specific to securities with an expiring contract.

This group is the identifying assumption of the study made testable. Without it every result
in Section~\ref{sec:results} is compatible with a settlement explanation and a calendar
explanation equally, and has to be reported with that caveat.

\input{generated/tables/universe}

\subsection{How the securities were chosen}

Membership is derived from the tape rather than taken from an index. An index membership list
downloaded today describes today: it includes securities listed since 2022, excludes
securities removed since, and carries present-day names for securities renamed in between.
Using one to select a 2022 sample silently mixes all three into the study.

The procedure is as follows.

\paragraph{Rank on control sessions only.} Turnover is averaged over the
\SelectControlSessions{} control sessions and the expiry sessions are excluded. This matters
more than it might appear. Turnover on an expiry Thursday is partly the phenomenon being
measured, so ranking securities on it would make the group variable a function of the
outcome, and part of any liquid-against-illiquid difference would be mechanical rather than
economic.

\paragraph{Require completeness.} A security must trade in every session in the sample. One
absent for a month cannot contribute that month's matched pair, and admitting it would make
each group's mean depend on which months its members happened to trade in.

\paragraph{Require enough activity to measure.} A security must clear a floor of
\SelectMinTrades{} trades in the median control session. This is a measurement criterion
rather than a size cut-off: below roughly this level a one-second book snapshot sees nothing
happen between consecutive observations, so the microstructure quantities are not identified
rather than merely small.

\paragraph{Require stable membership.} A security's turnover rank must not span more than
\SelectMaxRankSpan{} of the cross-section between its tenth and ninetieth percentile across
the year. One that is in the top decile in January and the bottom half in July is not
reliably in either group, and including it puts a security in one group that belonged in the
other for part of the sample.

\paragraph{Form the groups.} The \SelectGroupSize{} most traded and \SelectGroupSize{} least
traded of the surviving derivatives underlyings become the liquid and illiquid groups.

\paragraph{Match the placebo group.} Placebo securities are drawn from the surviving
securities with no derivative, by nearest-neighbour matching without replacement against the
illiquid group on two variables: turnover and price level, both in logarithms and
standardised so that neither dominates the distance.

Price level is not decoration. A quoted spread expressed in basis points depends on the tick
size relative to the price, so a placebo group of systematically cheaper shares would show
wider spreads for a mechanical reason with nothing to do with expiry, and the placebo
comparison would be biased before any data was examined. Matching without replacement matters
for a smaller reason: reusing one security as the match for several targets would make the
placebo group narrower than the group it stands in for.

\input{generated/tables/selection}

Table~\ref{tab:balance} reports how alike the groups ended up on the variables the matching
was meant to equalise. The placebo group's turnover range can only overlap the illiquid
group---a security traded as heavily as the liquid group while carrying no derivative is
close to a contradiction in terms, since heavy trading is roughly what derivatives
eligibility is granted for---so the placebo comparison speaks to the illiquid group and is
silent about the liquid one.

\input{generated/tables/balance}

\subsection{Derivatives eligibility}

Whether a security had an expiring contract cannot be read off the cash tape, so it is taken
from a published list of derivatives underlyings. Verified against a list for this run:
\SelectFoVerified{}. Where no list is available the two derivative groups are still formed on
turnover, no placebo group is produced, and the report states that expiry effects cannot be
separated from calendar effects.

\subsection{Window}

All measures are computed between \DesignWindowStart{} and \DesignWindowEnd{} IST. Where a
measure compares the window against the rest of the session it is expressed as a rate per
minute rather than as a total: the window is \DesignWindowMinutes{} minutes against roughly
three hundred and forty-five for the remainder, and a comparison of sums would reflect that
difference in length rather than any difference in behaviour.

Pre-open auction records share the session file with continuous trading and are excluded from
every return series. Otherwise the first continuous-session observation is differenced against
the auction price, which books the entire overnight move as one minute of intraday variance.
